\begin{figure*}[!ht]\centering
  \includegraphics[width=0.90\textwidth,height=0.78\textheight,keepaspectratio]{Fig3}
  \caption{\textbf{\emph{In-silico} reversal of the disease state.}
(\textbf{a}) The query signature is filtered on independent survival evidence: disease
direction against adjusted hazard ratio, with genes whose directions disagree set to zero.
(\textbf{b}) Positive control --- the official epiblast-to-nascent-mesoderm benchmark,
with the known regulators labelled and ranked (TBXT 1st of 8,450).
(\textbf{c}) Spatial copy-number score per domain, each domain compared with the remaining
spots of the same section so that all spots in a comparison share one library.
(\textbf{d}) Compound-level reversal: CMap rank against Cor-Spearman rank for the 1,826
compounds, with the 34 compounds common to all three scoring schemes highlighted, and their
mechanism composition.
(\textbf{e}) The candidate ranking is not modality-specific: a single-cell-derived and a
spatially derived query axis agree at Spearman $\rho=0.735$.
(\textbf{f}) A domain's marker stability scales with its sequencing depth
($\rho=-0.89$; D3 is the sole outlier and the only invasive-restricted domain).
(\textbf{g}) Of the 34 candidates, 26 fall in the top 100 for the ECM/interstitial domain
and 21 for the AT2 domain, against a chance expectation of 1.9.
(\textbf{h}) One library, three scores: the two correlation scores are near-redundant
($\rho=0.950$) whereas the CMap statistic shares only about a third of its ranking with
them.}\label{fig:3}
\end{figure*}
